\section{模型的评价与推广}

\subsection{主要结论}

对问题一，场景 A 的代价由依赖路径上的子图边界决定：每个边界带来一次 DDR 往返和 100 或 1000 周期的激活等待，并沿路径线性累积（命题~\ref{prop:path}）。据此设计的三种构造算法分别适合多分量、单一巨块和汇聚链结构，合法性由命题~\ref{prop:acyclic} 保证；逐用例选优与成组合并贡献了主要提升，核数补齐保证了 2--5 核结果单调。最终 5 核平均加速比为 \V{p1_k5_final_mean}，对含同步可达参考的达成率为 \V{p1_k5_achieve_mean}，剩余差距集中在 I 类小图上，受场景 A 固定延迟所限。

对问题二，同核子图合并为一个 Task 后，边界搬运由子图边界变为核边界（命题~\ref{prop:qb}），时间代价只随换核次数累积（式 \eqref{eq:path-B}），模型因此预测场景 B 适合细分，扫描数据与之一致。另一方面，跨子图的驻留数据加重了缓存压力（式 \eqref{eq:resident}），少数用例因此不如场景 A。最终 5 核平均加速比为 \V{p2_k5_final_mean}，I 类从问题一的 \V{p1_k5_famI_mean} 提高到 \V{p2_k5_famI_mean}。

对问题三，只读 L2 Cache 的命中同时受时序条件与容量条件约束（式 \eqref{eq:hit}），命中能否换来时间还取决于 DDR 是否紧张（式 \eqref{eq:ddr-busy}、\eqref{eq:gamma-bound}）。在官方配置附近，限制缓存收益的主要是读取时序而不是容量和带宽，锁步读取的对照实验提供了直接证据。两问共享候选池后，5 核只读 Cache 配置的平均加速比为 \V{p3_k5_cache_vs_single}，同核数缓存加速比为 \V{p3_k5_best_mean}，单例最高 \V{p3_k5_best_max}。

\subsection{模型的优点}

合法性由构造保证。三种构造算法都给子图一个随依赖严格递增的编号，生成的方案不需要事后修补就能通过评估程序的检查；后处理中可能成环的成组合并候选由评估程序检查后丢弃，不影响最终结果的合法性。

结果有参照，差距可以量化。第~\ref{sec:refs} 节的解析下界、零通信表调度和含同步可达参考逐步收紧，解析下界在 400 组（用例，核数）的实测中从未被突破。每个结果都能说明离可达水平还差多少，而不只是相对自身基线的提升。

模型能解释实验现象。命题~\ref{prop:path} 解释了场景 A 中为什么代价取决于路径深度，命题~\ref{prop:qb} 与式 \eqref{eq:path-B} 解释了场景 B 为什么适合细分，式 \eqref{eq:resident} 解释了少数用例场景 B 不如场景 A，式 \eqref{eq:hit} 与式 \eqref{eq:ddr-busy} 解释了只读 Cache 收益为什么小而集中。参数灵敏度中层带划分在两个场景下方向相反（表~\ref{tab:param}），也与模型的预测一致。全部数字来自官方评估程序的终版复评，1{,}400 个主结果逐位一致，三个问题都满足核数单调，未采用的路线也保留了可重现的实验记录。

\subsection{模型的不足}

简化模型只是松弛。式 \eqref{eq:obj-p1}、\eqref{eq:obj-p2} 忽略了 DDR 带宽争用和核内换出，它的最优解不一定是评估程序下的最优解，所以我们只用它指导算法设计和给出参照，最终结果依赖评估程序选优。含同步可达参考也只是贪心结果，4 核时问题一有 \V{p1_k4_above_ref} 个、问题二有 \V{p2_k4_above_ref} 个用例的实测值超过了它。

I 类中的小图仍是短板。5 核时 I 类的平均加速比，问题一为 \V{p1_k5_famI_mean}、问题二为 \V{p2_k5_famI_mean}，对可达参考的达成率为 \V{p1_k5_famI_achieve} 和 \V{p2_k5_famI_achieve}，低于 III 类（\V{p1_k5_famIII_achieve}、\V{p2_k5_famIII_achieve}）。对单核基准只有两万周期左右的小图，场景 A 的 1000 周期跨核等待使任何切分都得不偿失。

缓存的潜力没有充分挖掘，搜索成本也偏高。问题三 5 核的平均缓存加速比只有 \V{p3_k5_best_mean}，我们找到了限制收益的机理，但主动安排读取错位的两种尝试都没有成功。局部搜索用掉了 \V{ls_tries_total} 次官方评估，占全部 \V{evals_total} 次评估的 \V{ls_share_pct}，平均改进在几个百分点以内；更经济的做法可能是先用式 \eqref{eq:task-dur} 这类快速估计预筛候选，只把最有希望的交给评估程序。

\subsection{模型的推广}

三种构造算法的复杂度对核数是线性的，可以直接用于核数更多的 NPU。核数增加后，跨核同步在关键路径上的占比会继续上升，相位构造中的 $\lambda$ 和核数补齐会更重要。在实际编译器中，本文的结构分类与逐用例选优可以作为切图策略的选择器：先用 $\eta$、$\omega$ 等静态特征判断计算图类型，再决定用哪种构造算法，不必对每张图做完整搜索。

在缓存感知调度方面，式 \eqref{eq:hit} 给出了命中的时序条件。若能在调度阶段为每个共享张量指定领读核心，让跟读核心的 COPY\_IN 晚于领读完成时刻，就能把并发缺失转化为命中。第~\ref{sec:lockstep} 节的锁步实验表明，在部分用例上这样做的收益可能达到两成以上，这是值得继续研究的方向。
